\section{Conclusion and Limitations}
Restormer leads balanced restoration quality, and continued training explains most of the measured follow-up gain. The tested Nano cascade reduces detector AP and throughput while increasing input energy. Matched controls establish that resizing alone can remove a detection in one scene and that restoration at native resolution can help or hurt with the same detector.

Different initializations and one seed per model limit architectural comparisons. Six smoke backgrounds constrain the restoration benchmark. The labelled Nano and earlier local DAWN results use different geometries and scoring paths. The resolution diagnostic repeats five labelled objects, and camera predictions lack annotations. Longer board trials and a larger same-runtime, annotated resolution ablation would help quantify the causes and test generalization.
